\documentclass[10pt,a4paper]{amsart}
\usepackage[margin=19mm]{geometry}
\usepackage{amsmath,amssymb,mathtools}
\usepackage[scr=rsfs,cal=euler]{mathalfa}
\usepackage{xfrac}
\usepackage[unicode,hidelinks,bookmarks=false]{hyperref}
\newcommand{\ie}{i.e.\ }
\newcommand{\cf}{cf.\ }
\newcommand{\resp}{respectively\ }
\newcommand{\Subsection}{\subsection}
\providecommand{\nobreakdash}{\nobreak\hbox{-}}
\setlength{\emergencystretch}{2em}
\multlinegap0pt
\usepackage{amssymb}
\usepackage{enumerate}
\usepackage[shortlabels]{enumitem}
\newtheorem{lem}{Lemma}
\newcommand{\Cinf}{\mathbb{C}_{\infty}}
\DeclareMathOperator{\Ind}{Ind}
\DeclareMathOperator{\Res}{Res}
\newcommand{\bZ}{\mathbb{Z}}
\title{On some constancy of Hecke eigensystems for Drinfeld cuspforms of finite slope}
\author{Shin Hattori}
\date{Source: arXiv:2605.18016v2 (CC BY 4.0)}
\begin{document}
\maketitle

\noindent\emph{Source and licence.} On some constancy of Hecke eigensystems for Drinfeld cuspforms of finite slope. Version arXiv:2605.18016v2 (CC BY 4.0), distributed by arXiv under the Creative Commons Attribution 4.0 International licence, http://creativecommons.org/licenses/by/4.0/, as recorded in the arXiv licence metadata for this exact version and shown on the abstract page https://arxiv.org/abs/2605.18016v2. Full licence notice: \texttt{licenses/arxiv-2605.18016-CC-BY-4.0.txt}.\par\smallskip

\noindent\textit{Editorial scope.} Counted unit: the article's lem LemThetaFreeOfResolution (source0.tex lines 872--874). The article's complete original proof of it is delivered with it (source0.tex lines 875--913, one \texttt{proof} environment of 39 lines). No mathematics is written, repaired or re-derived: the statement and the proof are the article's own lines, in the article's own notation, and no retained line is substituted or re-flowed.  Retained uncounted as the source-local context the proof needs: lines 857--860.  Nothing was omitted from any retained span, so the package carries no omission record.  No retained line contains a citation, so the wrapper carries no bibliography and no bibliography entry.  No retained line contains a figure environment, an image-inclusion command or drawing code, so no image is delivered and the package declares no asset dependency.  Editorial formatting only, in addition: an AMS \texttt{amsart} preamble that restates the article's own declarations and macros used by the retained lines, namely the environments \texttt{lem} on separate counters, together with the standard packages those lines need.  The retained environments print in this document as Lemma 1 in order of appearance, whereas the article prints the same items with the numbering of its own full document.  The complete native source of the article as distributed in the arXiv e-print for this exact version is bundled unchanged as \texttt{sources/200832/source0.tex}.
\begin{equation}\label{EqnResolAreFree}
\tilde{L}_1=\bigoplus_{e\in \tilde{\Lambda}_{1}}[e] \bZ[\tilde{\Gamma}],\quad 
\tilde{L}_0=\bigoplus_{v\in \tilde{\Lambda}_{0}}[v] \bZ[\tilde{\Gamma}].
\end{equation}

\begin{lem}\label{LemThetaFreeOfResolution}
	For any $i\in\{0,1\}$, the right $\Cinf[\Theta_r]$-module $\tilde{L}_i\otimes_{\bZ[\Gamma]}V_k(\Cinf)$ is free of finite rank.
\end{lem}

\begin{proof}
	By (\ref{EqnResolAreFree}), we have an isomorphism of right $\Cinf[\Theta_r]$-modules
	\[
	\tilde{L}_i\otimes_{\bZ[\Gamma]}V_k(\Cinf)\to \bigoplus_{s\in \tilde{\Lambda}_{i}} \bZ[\tilde{\Gamma}]\otimes_{\bZ[\Gamma]}V_k(\Cinf).
	\]
	Thus we reduce ourselves to showing that the right $\Cinf[\Theta_r]$-module
	\[
	 \bZ[\tilde{\Gamma}]\otimes_{\bZ[\Gamma]}V_k(\Cinf)=\Ind_{\Gamma}^{\tilde{\Gamma}}\Res_{\Gamma}^{\tilde{\Gamma}}V_k(\Cinf)
	\]
	is free of finite rank.
	
	Let $\mathbf{1}$ be the trivial representation of $\Gamma$ over $\Cinf$. By the projection formula, we have a natural isomorphism of left $\Cinf[\tilde{\Gamma}]$-modules
	\[
	\pi:\Ind_{\Gamma}^{\tilde{\Gamma}}\Res_{\Gamma}^{\tilde{\Gamma}}V_k(\Cinf)\to (\Ind_{\Gamma}^{\tilde{\Gamma}}\mathbf{1})\otimes_{\Cinf} V_k(\Cinf)
	\]
	which is defined by
	\[
	\begin{aligned}
	 \bZ[\tilde{\Gamma}]\otimes_{\bZ[\Gamma]}V_k(\Cinf)&\to  (\bZ[\tilde{\Gamma}]\otimes_{\bZ[\Gamma]}\Cinf)\otimes_{\Cinf} V_k(\Cinf),\\
	 [\tilde{\gamma}]\otimes v&\mapsto [\tilde{\gamma}]\otimes 1 \otimes \tilde{\gamma}\circ v.
	 \end{aligned}
	\]
	On the target of the map $\pi$, we consider the trivial $\Theta_r$-action on $V_k(\Cinf)$ and
	the natural right $\Theta_r$-action on $\bZ[\tilde{\Gamma}]\otimes_{\bZ[\Gamma]}\Cinf$ given by
	\[
	([\tilde{\gamma}]\otimes 1)d=[\tilde{\gamma}\eta_{[1,d]}]\otimes 1
	\]
	for any $d\in \Theta_r$.
	Then the map $\pi$ is $\Theta_r$-equivariant. Moreover, with the latter $\Theta_r$-action, 
	we have an isomorphism of right $\Cinf[\Theta_r]$-modules
	\[
	\bZ[\tilde{\Gamma}]\otimes_{\bZ[\Gamma]}\Cinf\to \Cinf[\Theta_r],\quad [\eta_{[1,d]}]\otimes 1\mapsto [d].
	\]
	Hence we obtain an isomorphism of right $\Cinf[\Theta_r]$-modules
	\[
	\bZ[\tilde{\Gamma}]\otimes_{\bZ[\Gamma]}V_k(\Cinf)\simeq \Cinf[\Theta_r]^{\oplus k-1}.
	\]
	This concludes the proof.
\end{proof}

\end{document}
